\section{Null projections of the common point traces}
\label{sec:traces}

We now apply Theorem~\ref{thm:dust-sobolev} to a conformal equivalence
between circle domains.  The scalar functions used below are arranged to
be continuous even when a point component corresponds to a disk.

We first record a packing estimate closely related to
\cite[Lemma~2.1]{Ntalampekos2023}.  We include its short proof from the
Hardy--Littlewood Maximal Theorem \cite{Stein1970}.
\begin{lemma}[Ball synthesis]\label{lem:ball-synthesis}
Let $B_j\subset\R^2$ be balls of radii $d_j>0$.  Suppose there are
pairwise disjoint measurable sets $E_j\subset B_j$ with
$|E_j|\geq c d_j^2$, where $c>0$ is fixed.  For nonnegative
$(a_j)\in\ell^2$,
\[
 \left\|\sum_j\frac{a_j}{d_j}{\bf1}_{B_j}\right\|_{L^2(\R^2)}
 \leq C_c\left(\sum_j a_j^2\right)^{1/2}.
\]
\end{lemma}
\begin{proof}
Let $g\geq0$ belong to $L^2$, and let $Mg$ be its centered maximal
function.  For every $x\in B_j$, the average of $g$ on $B_j$ is at
most $4Mg(x)$, since $B_j\subset B(x,2d_j)$.  Consequently
\begin{align*}
 \int g\sum_j\frac{a_j}{d_j}{\bf1}_{B_j}
 &\leq 4\pi\sum_j a_jd_j\inf_{E_j}Mg\\
 &\leq C_c\left(\sum_j a_j^2\right)^{1/2}
       \left(\sum_j\int_{E_j}(Mg)^2\right)^{1/2}
 \leq C_c\|(a_j)\|_{\ell^2}\|g\|_2.
\end{align*}
Here the last inequality uses disjointness and the $L^2$ Maximal
Theorem.  Duality proves the assertion for finite sums, and monotone
convergence proves it for countable sums.
\end{proof}

For $e\in\{1,i\}$ write $\pi_e(w)=\Re(\overline e w)$.
\begin{proposition}[Common point traces]\label{prop:trace-null}
Let $\Omega,\Omega'$ be circle domains containing infinity, and let
$f:\Omega\to\Omega'$ be a conformal equivalence with $f(\infty)=\infty$.
Suppose $\partial\Omega$ is conformally removable.  Let $C$ be the set
of point components of $\Chat\setminus\Omega$ whose corresponding
components in $\Chat\setminus\Omega'$ are also points, and denote their
positions by $f_*(x)$, $x\in C$.  For each $e\in\{1,i\}$ and almost
every line $\ell$ parallel to $e$,
\[
 \bigl|\{\pi_e(f_*(x)):x\in C\cap\ell\}\bigr|=0.
\]
Here and below $|\cdot|$ for a subset of a line denotes Lebesgue length.
\end{proposition}
\begin{proof}
We fix $e$ throughout.  All complementary components on either side
lie in bounded planar regions.  Index those pairs $(K_j,K'_j)$ for
which at least one component is a disk of positive radius.  The index
set is finite or countable, since disjoint positive-radius disks have
disjoint nonempty interiors.  Write
\[
 K_j=\overline B(a_j,r_j),\qquad
 K'_j=\overline B(a'_j,s_j),
\]
allowing radius zero, and set $\beta_j=\pi_e(a'_j)$.  Disjointness and
boundedness give
\begin{equation}\label{eq:trace-radius-sums}
 \sum_jr_j^2<\infty,\qquad \sum_js_j^2<\infty.
\end{equation}
The other components are precisely the common points in $C$.

\smallskip\noindent
\emph{Scalar data and shrinking neighborhoods.}
Choose a smooth planar cutoff $\chi$, equal to one on a sufficiently
large disk containing the source complement and all the neighborhoods
introduced below, with its transition region compactly contained in
$\Omega$.  On $\Omega$ put
\[
 b=\chi\,\pi_e(f),
\]
and extend $b$ by $\beta_j$ on $K_j$, by $\pi_e(f_*(x))$ on $C$, and
by zero at infinity.  This is bounded.  Indeed, $f(\infty)=\infty$
implies that the image of a bounded source region is bounded.  The
area formula for the univalent map $f$ then shows that $b$ has finite
Dirichlet energy on $\Omega$; the cutoff-derivative term is supported
in a compact subset of $\Omega$.

Although $b$ need not be continuous at $K_j$, it is continuous at
every common point.  Moreover, the cluster values near $K_j$ lie in
$[\beta_j-s_j,\beta_j+s_j]$.  These statements include approaches
through other complementary components and follow from the
upper-semicontinuous correspondence in Lemma~\ref{lem:ends}.

Choose $\varepsilon_j>0$ tending to zero with
$\sum_j\varepsilon_j^2<\infty$, and put
\[
 I_j=(\beta_j-s_j-\varepsilon_j,\,
          \beta_j+s_j+\varepsilon_j),\qquad l_j=|I_j|.
\]
Thus $(l_j)\in\ell^2$.  Fix $A$ so large that all these intervals and
all values of $b$ lie in $(-A,A)$.  Choose positive $\Delta_j$ with
\begin{equation}\label{eq:trace-small-radii}
 \sum_{i>j}\Delta_i^2<\frac12\Delta_j^2;
\end{equation}
for example, a sufficiently small multiple of $2^{-j}$ suffices.
For each integer $m\geq1$, define the balloons
\[
 B_{j,m}=B(a_j,d_{j,m}),\qquad
 d_{j,m}=
 \begin{cases}
  (1+1/m)r_j,&r_j>0,\\
  \Delta_j/m,&r_j=0.
 \end{cases}
\]
These balls are all in one bounded region, and decrease to $K_j$ as
$m\to\infty$.  Choose radial Lipschitz functions
$0\leq\lambda_{j,m}\leq1$ that vanish on an open neighborhood of
$K_j$, equal one outside $B_{j,m}$, and satisfy
\begin{equation}\label{eq:trace-cutoff-gradient}
 |\nabla\lambda_{j,m}|
 \leq\frac{C_m}{d_{j,m}}{\bf1}_{B_{j,m}}
 \quad\hbox{almost everywhere}.
\end{equation}
One may take $C_m=2(m+1)$, interpolating radially between
$(1+1/(2m))r_j$ and $(1+1/m)r_j$ when $r_j>0$, and between
$d_{j,m}/2$ and $d_{j,m}$ otherwise.

For the moment fix $m$ and suppress it from the notation.  We will
transform the scalar value $b(z)$ before applying Sobolev extension.
Near $K_j$ the transformation should be constant on $I_j$, so it does
not distinguish the possible cluster values of $b$.  The cutoff
$\lambda_j$ imposes no restriction outside $B_j$.  Taking the
infimum of the cutoffs whose intervals contain a given scalar value
enforces these requirements simultaneously, including where the
intervals and neighborhoods overlap.  Integrating in the scalar
variable makes the spatial cost proportional to $l_j$, as the
estimate below will show.  Define
\begin{equation}\label{eq:trace-mask}
 L(z,t)=\inf_{j:\,t\in I_j}\lambda_j(z),\qquad
 F(z,t)=\int_{-A}^{t}L(z,s)\,\dd s,\qquad U(z)=F(z,b(z)),
\end{equation}
with the empty infimum equal to one.  We shall prove that $U$ is
continuous and belongs to $W^{1,2}(\Chat)$.

\smallskip\noindent
\emph{An energy density.}
Set
\begin{equation}\label{eq:trace-density}
 P=C_m\sum_j\frac{l_j}{d_j}{\bf1}_{B_j}.
\end{equation}
This belongs to $L^2(\C)$.  For the positive-radius sublist, the
interiors of the original disks $K_j$ are disjoint subsets of $B_j$
of area at least $\pi d_j^2/4$, so Lemma~\ref{lem:ball-synthesis}
applies.  For the zero-radius sublist use instead
\[
 E_j=B_j\setminus\bigcup_{i>j,\,r_i=0}B_i.
\]
These sets are pairwise disjoint: if $i<j$, then $E_i$ excludes
$B_j\supset E_j$.  Since all zero-radius balloons in this system
use the same factor $1/m$, \eqref{eq:trace-small-radii} gives
\[
 |E_j|\geq\pi d_j^2-\pi\sum_{i>j,\,r_i=0}d_i^2
 \geq\tfrac12\pi d_j^2.
\]
Apply Lemma~\ref{lem:ball-synthesis} separately to this sublist and
then add the two $L^2$ bounds.  Overlap between the two sublists is
irrelevant.

\smallskip\noindent
\emph{Continuity of the masks.}
We spell out the accumulation issue.  For fixed $m$, $d_j\to0$.
If distinct balloon supports accumulate at a common point $x$, their
centers and source components approach $x$.  By Lemma~\ref{lem:ends}
their target components approach $\{f_*(x)\}$; as
$\varepsilon_j\to0$, the corresponding intervals $I_j$ approach
the scalar $b(x)$.  Consequently, given $\delta>0$, there are a
neighborhood $V$ of $x$ and a finite initial list such that every
remaining balloon meeting $V$ has its interval in
$(b(x)-\delta,b(x)+\delta)$.  The finite-mask version of $F$ differs
from $F$ by at most $2\delta$ on $V$, uniformly in the upper endpoint.
Since the finite-mask version is continuous in its two variables and
$b$ is continuous at $x$, $U$ is continuous there.

For an indexed component $K_j$, choose a neighborhood $V_j$ of the
entire closed component on which $\lambda_j=0$ and $b\in I_j$.
This is possible because $I_j$ has a positive margin around all its
cluster values.  On $V_j$ the upper endpoint in
\eqref{eq:trace-mask} may be replaced by the fixed value $\beta_j$.
All sufficiently late intervals of balloons meeting a sufficiently
small neighborhood of $K_j$ lie in $I_j$.  To verify the uniformity
of this assertion, a failure would give $i_\nu\to\infty$ and
$z_\nu\in B_{i_\nu}$ with $\dist(z_\nu,K_j)\to0$, but
$I_{i_\nu}\not\subset I_j$.  Passing to a subsequence gives
$z_\nu\to x\in K_j$; the shrinking radii imply
$K_{i_\nu}\to\{x\}$.  Component correspondence puts their target
components arbitrarily close to $K'_j$, and the vanishing extra
margins then force $I_{i_\nu}\subset I_j$, a contradiction.
These late intervals have no effect on $L$ near $K_j$: there
$\lambda_j=0$, so $L(z,t)=0$ for $t\in I_j$.  Thus near $K_j$ the fixed-endpoint expression
is a finite continuous mask.  This proves continuity there, including
when $r_j=0<s_j$.

In $\Omega$ the cutoffs are locally finite, and near infinity they
are all one.  Hence the remaining continuity assertions are
immediate, with $U=A$ near infinity.

\smallskip\noindent
\emph{Sobolev regularity and the residual dust.}
Let $L_N,F_N,U_N$ denote the masks using just the first $N$ indices.
For $t\in[-A,A]$ and any $x,y$,
\begin{equation}\label{eq:trace-finite-variation}
 |F_N(x,t)-F_N(y,t)|
 \leq\sum_{j\leq N}l_j|\lambda_j(x)-\lambda_j(y)|.
\end{equation}
This follows from the elementary difference bound for finite minima,
integrated over the intervals on which a cutoff can occur.  Changing
the upper endpoint costs at most its change, since $0\leq L_N\leq1$.
It follows that $U_N$ has gradient bounded by $|\nabla b|+P$ locally
in $\Omega$.  On $V_j$, once $N\geq j$, its endpoint can be fixed
at $\beta_j$, and its gradient is bounded by $P$.

Bounded pointwise convergence of the finite masks gives local strong
$L^2$ convergence.  Weak Sobolev compactness therefore proves that
$U$ is locally Sobolev on
\[
 O=\Omega\cup\bigcup_jV_j,
\]
with the same gradient bounds.  These yield the finite total energy
estimate
\begin{equation}\label{eq:trace-off-dust-energy}
 \int_O|\nabla U|^2
 \leq2\int_\Omega|\nabla b|^2+2\int_\C P^2<\infty.
\end{equation}
Here the open union is integrated once; no multiplicity is introduced
by overlapping neighborhoods.

The residual set $B_* = \Chat\setminus O$ is compact and bounded,
and consists only of common point components.  Every connected
subset of $B_*$ lies in one component of
$\Chat\setminus\Omega$, hence is a singleton.  Thus $B_*$ is
totally disconnected.  Moreover $B_*\subset\partial\Omega$, and
so it is conformally removable by inclusion.  Theorem~\ref{thm:dust-sobolev},
applied to the continuous bounded function $U$ and
\eqref{eq:trace-off-dust-energy}, proves $U\in W^{1,2}(\Chat)$.
If $B_*$ is empty, this conclusion already follows from the local
Sobolev statement and the energy bound.  We have established the
conclusion for every mask system $m$; denote its function and density
by $U_m,P_m$.

\smallskip\noindent
\emph{A typical line and compression.}
Choose a line $\ell$ parallel to $e$ on which all $U_m$ are locally
absolutely continuous and all $P_m$ are locally integrable.
Almost every such line has these properties, since there are only
countably many systems.  We also require
\begin{equation}\label{eq:trace-line-tail}
 |\partial\Omega\cap\ell|=0,
 \qquad
 \sum_{j:\,B_{j,1}\cap\ell\ne\varnothing}l_j<\infty.
\end{equation}
The first condition holds on almost every line by
Lemma~\ref{lem:null-area} and Fubini.  For the second, write $\ell_\tau$
for the parallel line with transverse coordinate $\tau$.  Tonelli gives
\[
 \int_\R\sum_{j:\,B_{j,1}\cap\ell_\tau\ne\varnothing}l_j\,\dd\tau
 =2\sum_jl_jd_{j,1}<\infty,
\]
by Cauchy--Schwarz and \eqref{eq:trace-radius-sums}--\eqref{eq:trace-small-radii}.
Fix a line satisfying all these conditions.

For $N\geq1$ let
\[
 I^{(N)}=\bigcup_{\substack{j>N\\B_{j,1}\cap\ell\ne\varnothing}}I_j,
 \qquad
 H_N(t)=\int_{-A}^{t}{\bf1}_{\R\setminus I^{(N)}}(s)\,\dd s.
\]
Equation~\eqref{eq:trace-line-tail} implies $|I^{(N)}|\to0$.
For any $x\in C\cap\ell$, the distance from $x$ to each of the
first $N$ closed components is positive.  Some system $m$ therefore
has an open line neighborhood $J$ of $x$ missing its first $N$
balloons.  On $J$, every active cutoff belongs to the tail appearing
in $I^{(N)}$, and hence
\[
 L_m(z,t)\geq{\bf1}_{\R\setminus I^{(N)}}(t).
\]
The fixed-endpoint variation estimate extends to the infinite mask:
for any segment $[x,y]\subset\ell$ and fixed $t\in[-A,A]$,
\begin{equation}\label{eq:trace-segment-variation}
 |F_m(x,t)-F_m(y,t)|\leq\int_{[x,y]}P_m\,\dd s.
\end{equation}
Indeed \eqref{eq:trace-cutoff-gradient} and
\eqref{eq:trace-finite-variation} give this for finite masks, and
dominated convergence in the integration variable gives the limit.
The estimate uses only the globally Lipschitz cutoffs, so it remains
valid when the segment crosses indexed components where $b$ is
discontinuous.

If $x,y\in C\cap J$, compare the two upper endpoints at one fixed
location, and then use \eqref{eq:trace-segment-variation}.  The
nonnegative-integrand inequality above gives
\begin{align}
 |H_N(b(x))-H_N(b(y))|
 &\leq |U_m(x)-U_m(y)|+\int_{[x,y]}P_m\,\dd s\notag\\
 &\leq\int_{[x,y]}(|U_m'|+P_m)\,\dd s.
 \label{eq:trace-compression-bound}
\end{align}
For example, when $b(x)\leq b(y)$, first bound the compressed
increment by $F_m(x,b(y))-F_m(x,b(x))$ and then replace
$F_m(x,b(y))$ by $F_m(y,b(y))$.  The other order is identical.

The set $C\cap J$ has length zero.  A general elementary consequence
of \eqref{eq:trace-compression-bound} is therefore
\[
 |H_N(b(C\cap J))|=0.
\]
To see this with outer length, work first in a compact subinterval
of $J$.  Cover its null set of common points by an open set on which
the integral of $|U_m'|+P_m$ is arbitrarily small.  On each component
interval, \eqref{eq:trace-compression-bound} bounds the diameter of
the image of the common points by that integral.  Summing gives
arbitrarily small interval covers of the image.  A countable
exhaustion proves the assertion on $J$.

We recover the uncompressed values by the one-dimensional area
formula.  The common-point set is Borel, since it is the compact
complement with the countable union of the indexed closed components
removed.  Its trace map is continuous in the relative topology.
Thus $b(C\cap J)$ is analytic and Lebesgue measurable.  Put
$S=b(C\cap J)\setminus I^{(N)}$.  The nondecreasing Lipschitz
function $H_N$ has derivative one almost everywhere outside
$I^{(N)}$, and hence
\[
 |S|=\int_S|H_N'|
 =\int_\R\#\{t\in S:H_N(t)=u\}\,\dd u=0.
\]
The last integrand is supported on the null set
$H_N(b(C\cap J))$.  This reasoning allows both noninjectivity of
$H_N$ and its null set of derivative exceptions.

For fixed $N$, countably many of the neighborhoods $J$ suffice to
cover $C\cap\ell$, by the countable base of the line.  It follows
that $b(C\cap\ell)\setminus I^{(N)}$ is null.  Finally
$|I^{(N)}|\to0$, so $|b(C\cap\ell)|=0$.  Since $b(x)=\pi_e(f_*(x))$
on common points, this is the required conclusion.  Empty or finite
indexed lists are included by interpreting the corresponding sums
and tails in the evident way.
\end{proof}
